\section{Discussion}
\label{sec:discussion}

\subsection{What the $2\times$/$4.5\times$ number means operationally}

The headline numbers are unusual in the eVTOL/UAM literature because most
prior comparisons report incremental percentage-point improvements within a
single planning paradigm. The DREAM gap is multiplicative because it spans
two qualitatively different paradigms: \emph{configuration-blind} (B$_2$)
and \emph{configuration-aware} (B$_3$). The size of the gap quantifies the
operational cost of \emph{not} integrating runway-configuration awareness
into UAM corridor planning at an airport. Two thirds of LAX vertiport pairs
and seven eighths of SFO vertiport pairs that the static baseline would
declare infeasible are in fact reachable, on the same days, under the same
weather, with a 15-minute-resolution view of the active runway
configuration. Conversely, $\leq 20\,\%$ of path length is the price the
DREAM planner pays for that reachability on pairs the static baseline can
already serve. The net is unambiguous: airport-aware UAM planning needs the
geometric protection \emph{and} the operational awareness; neither alone
is sufficient.

\subsection{Why AUROC plateaus at $0.68$--$0.72$}

Across LR, RF, XGB, and MLP at both airports, the held-out AUROC clusters
within $\pm 0.05$ of $0.70$ (\cref{tab:risk}). This is consistent across
model classes, which means the ceiling is in the features rather than in the
model. Our current feature set is dominated by static geometric distances
(d$_\mathrm{OLS}$, d$_{\mathrm{rwy}}$, d$_{\mathrm{app}}$, d$_{\mathrm{dep}}$)
and aggregated meteorological state; neither captures the
\emph{per-aircraft kinematic} state that TCAS-style closest-point-of-approach
predictors use~\citep{rtca_dolca,lefebvre2018}. Engineering CPA features --
relative speed and bearing of each ADS-B track at the candidate
midpoint -- would, by analogy with the ATM conflict-prediction literature,
push AUROC into the $0.80$--$0.90$ band. We flag this as the most
impactful immediate extension.

\subsection{Why calibration succeeds where discrimination is bounded}

Split-conformal prediction \citep{vovk_conformal,angelopoulos2021gentle}
guarantees coverage $1 - \alpha$ on exchangeable data, regardless of
whether the underlying classifier discriminates well. The SFO XGB model
achieves coverage $0.912$ with AUROC $0.716$ -- i.e.\ the
prediction sets are correctly sized, even though the point predictions
are noisier than the experiment plan asked for. For an operational
safety-critical context, the right level of claim to make on the current
feature set is calibration, not discrimination. We have done so throughout.

\subsection{Vertiport siting (H3)}

The per-pair heatmaps (\cref{fig:vertiport-heatmap}) replicate at both
airports the qualitative pattern that landside-edge vertiports ($V_2$)
dominate terminal-rooftop vertiports ($V_3$) on joint feasibility -- the
proposal's H3. The rooftop is geometrically the tightest case under
both-sides-closed Annex~14 protection because its OFV funnel sits between
the parallel runway corridors; the off-fence $V_1$ is feasible mainly
because it lies outside the inner-horizontal surface; the city-end $V_4$
(Union Station for LAX, Salesforce Transit Center for SFO) is sometimes
infeasible because of overflight constraints on long city-side corridors.
The operational reading is that landside-edge intermodal facilities such as
the LAX Intermodal Transportation Facility West are the right primary
deployment target for eVTOL airport shuttles in this generation, with
rooftop vertiports reserved for follow-up deployments that have access to
finer planning resolution or active-runway-corridor coordination.

\subsection{Threats to validity and limitations}
\label{sec:discussion-limitations}

We list the threats explicitly:
\begin{itemize}
\item \textbf{No real eVTOL operational data.} All eVTOL trajectories are
    counterfactually injected into real fixed-wing operational scenes.
    Conformal coverage is therefore a calibration claim about the
    \emph{injected} distribution, not about a deployed service.
\item \textbf{Coverage variance in adsb.lol.} The volunteer-receiver
    network shows a $\sim 50\,\%$ drop in coverage on 2024-08-{16, 23, 30}
    at both airports (\cref{tab:data}). The $B_3 \supsetneq B_2$ ranking
    is robust to this drop; absolute risk-model AUROC would benefit from
    a denser feed (FAA SWIM SWIFT).
\item \textbf{Annex 14 parameters are public-secondary-source.} The values
    in \cref{tab:ols-params} are the publicly documented Code~4 precision
    approach values, not the ICAO normative text. Per-airport calibration
    requires a YAML override.
\item \textbf{Risk-field AUROC plateau.} Discussed in
    \cref{sec:discussion}. The model-class invariance localises the gap to
    the feature set.
\item \textbf{Memory-bound corridor-closure-rate KPI.} Loading
    five days of envelope tensors simultaneously exceeded our 54~GB
    Featurize box and was disabled for this report; the streaming
    variant is a $\sim 1$-day refactor.
\item \textbf{B$_4$ risk-field-cost sweep.} The per-cell risk grid export
    was deferred in favour of the headline B$_2$/B$_3$ comparison. The
    code path is in place; we leave the B$_4$ sweep as immediate future
    work.
\item \textbf{UTC-aligned peak hours.} The three peak hours of our sweep
    are UTC 08, 11, and 17. The LAWA-defined ground-traffic peaks are
    local PT, which corresponds to a 7-hour shift in UTC. We retain UTC
    hours throughout for cross-airport consistency, at the cost of not
    perfectly aligning with the LAWA local peaks. A future revision should
    rerun with airport-local peak hours and bootstrap confidence intervals
    on the $B_3 \to B_2$ difference per pair.
\end{itemize}

\subsection{Regulatory implications}

FAA Engineering Brief 105A~\citep{faa_eb_105a} explicitly requires that
vertiport approach and departure paths be independent of the active runway
paths \emph{where practicable}. DREAM is, to our knowledge, the first
quantitative tool that says how often this independence requirement is in
fact \emph{practicable} on a per-vertiport basis at a real airport. The
operational implication is that the practicability of a vertiport siting
should be assessed on the \emph{dynamic} envelope $\Edyn$ rather than on the
static $\Astatic$; otherwise the safety case is over-conservative.

\subsection{Future work}

The immediate extensions are:
(i)~\emph{TCAS-CPA features} -- per-aircraft kinematic features should
break the AUROC plateau;
(ii)~\emph{full B$_4$ sweep} with the risk-grid term active;
(iii)~\emph{streaming corridor-closure-rate KPI} -- a memory-bounded
re-implementation of the safety-axis aggregator;
(iv)~\emph{SWIM SWIFT migration} -- moving from \texttt{adsb.lol} to FAA
SWIM SWIFT for production-grade ADS-B with no coverage gaps;
(v)~\emph{additional hub airports} -- Boston Logan, Chicago O'Hare,
Hartsfield-Jackson Atlanta would round out the cross-airport replication
across different runway-geometry classes;
(vi)~\emph{integration with vertiport scheduling}~\citep{vascik2019} --
the corridor-feasibility table from DREAM is the natural input to a
vertiport-capacity envelope, closing the gap to the existing scheduling
literature.
